\section{Open items}
\label{sec:open}

\textbf{Restarted 2026-07-20.} This list previously mixed genuine open
task-graph/algorithm work with validation-run tracking (results
against Hu et al.\ 2017/Smith et al.\ 2021/STARBENCH, and the e10
diagnostic history). Validation status now lives exclusively in
Section~\ref{sec:validation} (its own summary table,
Table~\ref{tab:validation-targets}); every settled bug fix or shipped
feature that had accumulated here was moved to
Section~\ref{sec:bugs} or the relevant part of Section~\ref{sec:algorithm}/
\ref{sec:taskgraph}. This list is renumbered from T1 and contains only
currently-open task-graph/algorithm items -- nothing here tracks
validation-against-a-paper status.

\begin{itemize}
  \item[$\square$] \textbf{T1} -- Confirm or rule out the
    \code{radiation\_level} orphaned-link risk
    (Section~\ref{sec:not-regular}) on the clump examples specifically.
    It requires two neighbouring regions' \code{hydro.super} to land
    at different tree depths, which uniform-density ICs never trigger
    but a $100\times$ density clump plausibly could. Rebuild with
    \code{--enable-debugging-checks} and run
    \code{StromgrenSphereClumpSmith2021} at a resolution fine enough
    to actually split the tree near the clump boundary; the
    \code{error()} at \code{cell.c:1387} will fire if it's hit.

  \item[$\square$] \textbf{T2} -- MPI is explicitly not supported yet
    for the radiation loop (\texttt{error("MPI is not yet
    implemented")} in both \code{task.c} and \code{cell\_unskip.c}
    for \texttt{nr\_nodes > 1}). Not a bug, just an explicit gap to be
    aware of before attempting a multi-node run with this feedback
    channel active.

  \item[$\square$] \textbf{T3} -- The e10 diagnostic
    (Section~\ref{sec:e10} -- an algorithm-behaviour check, not a
    validation target, see the note there) has an unexplained residual:
    even with the search-radius-expansion cap loosened and the
    adaptive rebuild cadence on, the simulated D-type growth curve
    still shows a gap to the analytic solution that neither fix
    closes. Both identified contributing factors (the cap, the fixed
    cadence) are confirmed real and exhausted as candidate root causes
    at their current settings -- the residual points at something else,
    not yet identified. A $T_i=10^4\,$K variant of the same combined
    setup passed the check outright (Section~\ref{sec:e10}), but three
    variables changed at once relative to the original mismatch, so
    $T_i$ is not cleanly isolated as the explanation. Not attempted:
    instrumenting which retry branch (buffer-full vs.\ radius-expansion)
    dominates once the cap/cadence fixes are both in place.

  \item[$\square$] \textbf{T4} -- User's stated direction (2026-07-19):
    the adaptive rebuild cadence (Section~\ref{sec:adaptive-cadence})
    is intended to become the \emph{default}, not stay opt-in
    indefinitely. \textbf{Deliberately not done yet} (2026-07-20): the
    user wants more testing first, so
    \code{HII\_adaptive\_rebuild\_cadence} stays off by default.

  \item[$\square$] \textbf{T5} -- Pass real ionizing/heating rates
    into Grackle's chemistry solve (Grackle's own \code{RT\_heating\_rate}/
    \code{RT\_HI\_ionization\_rate} fields), instead of this model's
    instant flag-and-floor scheme. User's explicit position
    (2026-07-19): worth testing, \emph{not} for the current STARBENCH
    validation effort (both \citet{Hu2017} and \citet{Smith2021}
    themselves use the same flag-and-floor approach this model does, so
    rate-coupling would solve a different problem than the one
    currently being validated against), and not before it -- a
    standalone follow-up. See memory
    (\code{reference\_gizmo\_fire\_rt\_rate\_coupling.md}) for how
    FIRE-3 actually does this (self-consistent solve, Stromgren-value
    incident flux, self-shielding off) and why it needs an actual local
    radiation field, not just a number handed to Grackle. \textbf{Update
    (2026-07-20)}: confirmed Grackle also has an equally-real, separate
    coupling point specifically for H$_2$ photodissociation
    (\code{RT\_H2\_dissociation\_rate}, gated on
    \code{primordial\_chemistry}\,$\in\{2,3\}$, i.e.\ this project's
    modes 2/3) and for photoelectric heating
    (\code{use\_isrf\_field}/\code{isrf\_habing}, a genuinely per-particle
    field, not just the scalar \code{interstellar\_radiation\_field}) --
    see Section~\ref{sec:fuv-isrf} for the full trace into Grackle's own
    source. Nothing needs changing inside Grackle for either channel;
    the entire gap is T6's local FUV estimate below.

  \item[$\square$] \textbf{T6} -- Variable ISRF (interstellar radiation
    field), user's design idea (2026-07-19), not yet implemented.
    Motivation: proper Grackle coupling (T5) needs an ISRF estimate, and
    the common approach (reusing the gravity tree for this) is, in the
    user's own words, ``not optimal nor ideal'' and would require new
    tasks for a lot of engineering cost. \textbf{Proposed alternative,
    cheaper and (in the user's assessment) ``not less correct''}: inject
    FUV radiation onto gas particles within the source's kernel, then
    propagate the FUV field by smoothing over gas neighbours -- possibly
    using the analytical solution to a hyperbolic diffusion equation,
    explicitly \emph{not} an M1-closure scheme, specifically to avoid
    the $\text{CFL}\sim1/c$ constraint M1 would impose. Not scoped or
    attempted. \textbf{Update (2026-07-20)}, Section~\ref{sec:fuv-isrf}:
    photoelectric heating and H$_2$ (Lyman-Werner) dissociation are
    \emph{not} the same photon band (Habing/PE $\sim6$--$13.6\,$eV vs.\
    LW $11.2$--$13.6\,$eV, a hard sub-band). \textbf{Decided}: rather
    than assume a fixed ISRF spectral shape to split one proxy in two,
    track two explicit band luminosities per star, $L_{\rm FUV}$
    ($6$--$11.2\,$eV) and $L_{\rm LW}$ ($11.2$--$13.6\,$eV) --
    Section~\ref{sec:fuv-two-bands} -- summed for
    \code{isrf\_habing}/photoelectric heating, with $L_{\rm LW}$ alone
    feeding \code{RT\_H2\_dissociation\_rate}. \textbf{Also decided}:
    injecting these onto kernel-neighbour gas
    (Section~\ref{sec:fuv-propagation}) is not itself a smooth ISRF
    (sharp kernel cutoff, resolution-dependent reach); propagation
    should solve a \emph{screened} (Helmholtz-type) diffusion equation,
    not a plain one -- its point-source solution is a Yukawa
    $e^{-r/\lambda}/r$, matching real geometric-dilution-plus-dust-
    extinction falloff, unlike plain diffusion's unscreened,
    unphysical $1/r$. The CFL win comes from solving this
    \emph{steady-state} each update (light-crossing time
    $\ll$ hydro/cooling timescales), not from "diffusion" as such -- a
    time-marched parabolic diffusion PDE has its own, potentially just
    as restrictive, $\Delta t\lesssim\Delta x^2/D$ constraint. Still not
    scoped: how to solve that steady equation inside SWIFT's
    SPH+task-graph architecture without a full global linear solve
    every step.

  \item[$\square$] \textbf{T7} -- Extend \texttt{pychem}
    (\texttt{\textasciitilde/programs/pychem/}, already the source of
    this project's GEAR chemical-yield tables) to emit proper
    $(M,Z)$-dependent stellar radius/luminosity/effective-temperature
    (and hence ionizing photon rate) tables, replacing the current
    single-metallicity, mutually-disagreeing piecewise power-law fits
    in \code{radiation.c} (Section~\ref{sec:blackbody} documents the
    disagreement in detail). A standalone project, not a small in-place
    fix. When it lands, delete the dead
    \code{radiation\_get\_individual\_star\_temperature}
    (Section~\ref{sec:blackbody}). \textbf{Scope widened (2026-07-20)}:
    $L_{\rm FUV}$ and $L_{\rm LW}$ (Section~\ref{sec:fuv-two-bands}) are
    two more band integrals of the same underlying stellar-atmosphere
    model this project would tabulate -- not a separate
    table-generation effort from $Q_H$/$L_{\rm bol}$, just two more
    threshold pairs on the same curve.

  \item[$\square$] \textbf{T8 -- citation partially resolved
    (2026-07-20), validation gap still open} -- The radiation-pressure
    momentum recipe (Eq.~\ref{eq:rp-momentum-rate}--\ref{eq:rp-tauIR},
    Section~\ref{sec:radiation-pressure}: single UV scattering plus
    $(1+\tau_{\rm IR})$ dust-IR-trapping boost, $\kappa_{\rm IR}=10\,{\rm
    cm^2\,g^{-1}}$ at $Z_\odot$) is not unique to one paper -- it is a
    widely shared recipe across \citet{Hopkins2018} (FIRE-2, already
    cited elsewhere in this logbook), \citet{Agertz2013}, and
    \citet{Marinacci2019} (the SMUGGLE model), and the $\kappa_{\rm
    IR}=10\,{\rm cm^2\,g^{-1}}\,(Z/Z_\odot)$ normalization matches a
    commonly-quoted value across this literature. \textbf{Not
    resolved}: which single paper (if any) this codebase's exact
    implementation traces to -- full-text access to
    \citet{Agertz2013}/\citet{Marinacci2019}'s own equations was not
    available in this pass to confirm their precise $\kappa_{\rm IR}$
    value against this code's. \textbf{Separately, still fully open}:
    pinning a citation does not validate the physics -- no test in
    this codebase exercises the radiation-pressure channel at all
    (Section~\ref{sec:rp-open}), so the numerical fix in
    Section~\ref{sec:rp-hit-flag-bug} (momentum now actually applied)
    has not been checked against any analytic expectation.
    \textbf{Proposed two-tier plan (2026-07-20)},
    Section~\ref{sec:rp-validation-plan}: (1) a cheap
    momentum-conservation self-check ($\int \dot p_{\rm rad}\,dt$ from
    the code's own logged $L_{\rm bol}$/$\tau_{\rm IR}$ vs.\ actual gas
    momentum gained -- code correctness only, not a physics check); (2)
    a genuine fourth validation target against
    \citet{KrumholzMatzner2009}'s analytic radiation-pressure-driven
    shell radius(t) solution, the direct radiation-pressure analog of
    the Spitzer D-type solution STARBENCH validates for photoionization,
    with \citet{Draine2011}'s static-equilibrium dusty-H\,\textsc{ii}
    solutions as a possible cross-check. Neither tier implemented yet;
    tier~1 is queued first (cheap, would have caught the
    \code{hit\_by\_radiation} bug on its own), tier~2 is not yet scoped
    as its own example directory.

  \item[$\square$] \textbf{T9} -- \code{space\_regrid}'s top-level
    cell grid never re-refines once coarsened, even after the
    \code{h\_hii} that coarsened it shrinks back down (e.g.\ a star
    dies). Traced in code (2026-07-20), not yet fixed.
    \textbf{Per-star and per-cell bookkeeping is already correct}:
    \code{feedback\_common.c:372-378} explicitly resets
    \code{sp->h\_hii = 0.0} when a star dies or ages out of HII
    eligibility (comment: ``This prevents dead particles with large
    h\_hii to force the tasks at higher levels than necessary'',
    final value preserved in
    \code{sp->tracers\_data.final\_HII\_radius}), and
    \code{c->stars.h\_hii\_max} (the per-cell aggregate
    \code{space\_regrid} reads, Section~\ref{sec:space-regrid-hhii}) is
    recomputed fresh from current \code{sparts[i].h\_hii} at every
    tree rebuild/drift (\code{cell.c}, \code{cell\_drift.c},
    \code{space\_split.c}) -- so it does shrink correctly once a
    star's \code{h\_hii} resets. \textbf{The gap is one level up}:
    \code{space\_regrid.c}'s own regrid trigger,
    \begin{quote}
    \code{if (s->cells\_top == NULL || cdim[0] < s->cdim[0] || ...)}
    \end{quote}
    only fires when the newly-computed \code{cdim} is \emph{smaller}
    (coarser) than the grid's current \code{cdim} -- never when it
    would come out larger (finer) again. This is a general,
    pre-existing SWIFT property, true for ordinary hydro \code{h\_max}
    too, not introduced by this branch: once cells coarsen, they stay
    coarse for the rest of that run, regardless of why \code{h\_max}
    grew. \code{h\_hii} is the first mechanism in this codebase where a
    \emph{single, brief, transient} event (one star's HII region,
    possibly capped only by the periodic floor of 3,
    Section~\ref{sec:space-regrid-hhii}) can trigger this, unlike ordinary
    hydro resolution needs which are usually sustained, not transient.
    \textbf{Practical concern}: in a real multi-star simulation, one
    star's HII region growing large enough to coarsen \code{cdim} --
    even briefly, even down to the periodic floor -- would coarsen the
    whole box for the remainder of that run, even after that star dies
    and every other star stays small. Not fixed here: needs a design
    decision (e.g.\ periodically re-attempt a finer regrid when
    \code{h\_max} has dropped and stayed dropped for some margin, or
    track \code{h\_hii}'s contribution to \code{h\_max} separately so
    it cannot ratchet the whole grid the way sustained hydro resolution
    needs legitimately should).

  \item[$\square$] \textbf{T10 -- confirmed (2026-07-21), decision
    still open} -- STARBENCH re-validation
    (Section~\ref{sec:starbench-current}): the plain committed default
    build \textbf{FAILS at 22.76\%} (tolerance 10\%); building with
    \code{IONIZATION\_FEEDBACK\_DEBUG\_FIXED\_IONIZED\_TEMPERATURE\_K=1e4}
    on top of the exact same (unchanged) committed \code{run.sh}/
    \code{params.yml} \textbf{PASSES at 6.12\%}. $T_i$ convention is
    confirmed as the sole differing factor -- this example's own
    search-radius cap and rebuild cadence have not changed since first
    committed (checked via \code{git log --follow}, correcting an
    earlier imprecise statement that they had), and the search-radius
    expansion/retry mechanism was directly verified not to bind in the
    passing run (\code{HIIRegionRadii(t)} grows smoothly to its natural
    $2.46\,$pc equilibrium with monotonically decelerating growth
    factors, never approaching the $4.9\,$pc ceiling). \textbf{Still
    open}: deciding whether matching STARBENCH's flat-$T_i$ convention
    should become an actual default-buildable option (not just a debug
    flag), or whether this code's natural non-isothermal $T_i$ is the
    intended, more physical behaviour and STARBENCH comparison is
    understood to always require this explicit build -- the committed
    example's README now documents the compile flag either way, so
    this is a design decision, not a reproducibility gap. Per the
    restart ordering (Section~\ref{sec:validation-targets}), STARBENCH
    can now be considered validated (with the documented build), so
    the next target is Hu et al.\ (2017)/Smith et al.\ (2021) 4-star.
\end{itemize}
